%=============================================================================
% Chapter 6: Frontend Implementation
%=============================================================================
\chapter{Frontend Implementation}

The FloodEye frontend is a Next.js 16 App Router application. All pages live under either the root \texttt{/} (landing) or the nested \texttt{/dashboard/} route group. This chapter describes each page in detail.

\section{Landing Page (\texttt{app/page.tsx})}

\subsection{Purpose}
The landing page serves as the public-facing introduction to FloodEye. It is designed to be visually striking and informative, communicating the project's purpose to new visitors before they enter the monitoring dashboard.

\subsection{Components}
\begin{itemize}[leftmargin=1.5em]
  \item \textbf{LandingNav:} Fixed top navigation bar with logo, navigation links, and a ``Go to Dashboard'' call-to-action button. Includes blur-glass styling.
  \item \textbf{BackgroundSlideshow:} An auto-advancing slideshow of flood monitoring imagery rendered as a full-viewport background layer.
  \item \textbf{ScrollytellingSections:} The primary content component. Uses GSAP ScrollTrigger to animate content sections as the user scrolls: Hero, Problem Statement, Features Showcase, Technology Stack, Team Profiles (with individual member photos from \texttt{/public/}), and a call-to-action.
  \item \textbf{AssamNodeMap:} A landing-page specific map component showing the geographic coverage area of the monitoring network in Assam.
  \item \textbf{PWAInstallPrompt:} A floating prompt that appears when the browser's beforeinstallprompt event fires, inviting the user to install the PWA.
\end{itemize}

\subsection{Key Features}
\begin{itemize}[leftmargin=1.5em]
  \item Lenis smooth scroll integration for physics-based scroll inertia.
  \item GSAP \texttt{fromTo} and \texttt{ScrollTrigger} animations for each section entrance.
  \item Fully responsive layout adapting from mobile through desktop breakpoints.
  \item Team member profile cards with photos from the public directory.
\end{itemize}

\section{Dashboard Home (\texttt{app/dashboard/page.tsx})}

\subsection{Purpose}
The main dashboard provides a comprehensive real-time overview of all sensor readings, device status, live charts, alerts, and the flood risk prediction in a single scrollable view.

\subsection{Layout Structure}

The dashboard is organised in four rows:

\begin{enumerate}[leftmargin=1.5em]
  \item \textbf{Row 1:} Five sensor cards (Temperature, Humidity, Pressure, Altitude, Water Level) in a responsive grid (1 col on mobile \textrightarrow{} 5 cols on XL screens).
  \item \textbf{Row 2:} A 2:1 grid containing the Water Level Trend chart (left, spanning 2 columns) and a right column with ComfortScoreCard, FloodPredictionCard, and DeviceStatusCard stacked vertically.
  \item \textbf{Row 3:} The AlertPanel spanning the full width.
  \item \textbf{Row 4:} The interactive sensor location map (MapLibre GL, loaded dynamically with SSR disabled).
\end{enumerate}

\subsection{Sensor Cards (SensorCard Component)}

Each SensorCard renders:
\begin{itemize}[leftmargin=1.5em]
  \item Current reading with unit.
  \item Trend arrow icon (up / down / stable) calculated by comparing the current reading against the previous one.
  \item Delta value showing the numeric change from the previous reading.
  \item A sparkline mini-chart rendered using the \texttt{historyData} prop.
  \item A status badge (normal / warning / critical) with colour coding.
  \item Offline-state: Cards are labelled ``(Cached)'' when the device is offline.
\end{itemize}

Alert thresholds per card:
\begin{itemize}[leftmargin=1.5em]
  \item Temperature $>$ 35°C: warning status.
  \item Humidity $>$ 80\%: warning status.
  \item Water Level $<$ 20 cm: critical; $<$ 40 cm: warning.
\end{itemize}

\subsection{Water Level Trend Chart}
An interactive Recharts \texttt{CustomLineChart} displays the last N data points of water level (cm) vs.\ time. A ``Live'' or ``Cached'' pill badge indicates the current data freshness.

\subsection{Comfort Score Card}
Calculates and displays the Environmental Comfort Score as a percentage derived from temperature and humidity. Score ranges map to qualitative labels (Excellent / Good / Moderate / Poor).

\subsection{Flood Prediction Card}
A compact summary of the ML prediction result: predicted water level in 4 hours, risk level badge (Low / Moderate / High / Critical), and flood probability percentage.

\subsection{Device Status Card}
Displays ESP32 online/offline status, ThingSpeak connection status, and last upload time derived from the \texttt{DeviceStatus} object.

\subsection{Alert Panel}
The AlertPanel component renders a horizontally scrollable list of active alerts, with severity-coded colour bands. It links to the full Alerts page.

\subsection{Offline Behaviour}
When the ESP32 is unreachable:
\begin{itemize}[leftmargin=1.5em]
  \item An orange \textbf{OfflineBanner} is shown at the top with the time of the last live reading.
  \item Sensor cards show ``(Cached)'' labels.
  \item The Live badge on the chart changes to ``Cached''.
  \item A \textbf{NoDataState} component replaces the dashboard if no cached data exists at all.
\end{itemize}

\section{Analytics Page (\texttt{app/dashboard/analytics/page.tsx})}

\subsection{Purpose}
Provides individual time-series charts for all five sensor channels, each rendered in its own card.

\subsection{Components}
\begin{itemize}[leftmargin=1.5em]
  \item \textbf{TimeFilter:} A pill-tab selector for 1h, 6h, 24h, 7d, 30d windows. Selecting a window triggers a new \texttt{/api/history?range=\{window\}} request.
  \item \textbf{Five CustomLineChart cards:} Temperature (°C), Humidity (\%), Atmospheric Pressure (hPa), Altitude (m), Water Level (cm). The Water Level card uses a dark gradient background to visually distinguish it as the primary safety metric.
  \item The Pressure chart spans two columns (\texttt{lg:col-span-2}) as a wide chart for better readability.
\end{itemize}

\section{History Page (\texttt{app/dashboard/history/page.tsx})}

\subsection{Purpose}
Displays statistical summaries of sensor history and provides data export functionality.

\subsection{Components}
\begin{itemize}[leftmargin=1.5em]
  \item \textbf{StatCard grid:} Six cards showing average, min, and max for temperature, humidity, and water level across the current history window.
  \item \textbf{Data table:} A scrollable table of the last 100 readings with columns: Time, Temperature, Humidity, Pressure, Altitude, Water Level.
  \item \textbf{Export buttons:}
    \begin{itemize}
      \item \textbf{CSV Export:} Generates a Blob from the history array and downloads it as \texttt{floodeye\_logs\_\{date\}.csv}.
      \item \textbf{PDF Export:} Uses jsPDF + jspdf-autotable to generate a formatted PDF with the FloodEye header and a data table.
    \end{itemize}
\end{itemize}

\section{Predictions Page (\texttt{app/dashboard/predictions/page.tsx})}

\subsection{Purpose}
A dedicated full-page view of the ML flood forecasting results with admin diagnostic tools.

\subsection{Components}
\begin{itemize}[leftmargin=1.5em]
  \item \textbf{FloodPredictionPanel:} A large card displaying the predicted water level, a circular risk gauge, flood probability bar, confidence score, feature contributions list (Water Level, Temperature, Humidity, Pressure), and a timestamp of the last inference.
  \item \textbf{Admin Diagnostics (admin role only):}
    \begin{itemize}
      \item Model Architecture card: Type (CNN-BiLSTM-Att), Input Shape [48, 11], Output (4hr Forecast), MAE (Train: 3.42 cm).
      \item Inference Engine card: Provider (TensorFlow.js/WebGL), live status (COMPUTING/READY), last run timestamp.
      \item Settings CTA card: Links to the Settings page.
    \end{itemize}
\end{itemize}

\section{Alerts Page (\texttt{app/dashboard/alerts/page.tsx})}

\subsection{Purpose}
The full alert log with severity filtering and per-alert dismissal.

\subsection{Components}
\begin{itemize}[leftmargin=1.5em]
  \item Summary badges showing counts of Critical, Warning, and Info alerts.
  \item \textbf{AlertCard:} Each alert card displays: severity icon, type label, severity badge, timestamp (relative, e.g.\ ``3m ago''), message text, and a dismiss button.
  \item Alerts are ordered critical-first, then warning, then info.
  \item ``Clear All'' button clears the log and resets all alert cooldown timers.
\end{itemize}

\textbf{Alert types implemented:}
\begin{itemize}[leftmargin=1.5em]
  \item HIGH\_TEMPERATURE, HIGH\_HUMIDITY, RAPID\_PRESSURE\_CHANGE
  \item HIGH\_WATER\_LEVEL, RAPID\_WATER\_RISE, OBJECT\_TOO\_CLOSE, SENSOR\_OFFLINE
\end{itemize}

\section{Settings Page (\texttt{app/dashboard/settings/page.tsx})}

\subsection{Purpose}
Allows admin users to configure alert thresholds, refresh interval, and simulation mode.

\subsection{Access Control}
The Settings page checks \texttt{userRole === "admin"} from the TelemetryProvider context. Non-admin users see an access-restricted screen with a ShieldAlert icon.

\subsection{Settings Sections}
\begin{itemize}[leftmargin=1.5em]
  \item \textbf{Connection Settings:} Refresh interval selector (5s, 15s, 30s, 60s). Simulation mode toggle (uses synthetic data instead of live ThingSpeak readings).
  \item \textbf{Alert Thresholds:} Input fields for Temperature threshold (°C), Humidity threshold (\%), Water Level threshold (cm), Pressure threshold (hPa).
  \item \textbf{Save:} Persists settings to the TelemetryProvider context (in-memory; not yet persisted to a database in the current implementation).
\end{itemize}

\section{Devices Page (\texttt{app/dashboard/devices/})}

The Devices section provides device management views showing registered ESP32 nodes and their connection status. In the current implementation, this page displays the status of the single configured ThingSpeak channel.

\section{Responsive Design}

All dashboard pages are fully responsive using Tailwind CSS breakpoint utilities:
\begin{itemize}[leftmargin=1.5em]
  \item \textbf{Mobile (default):} Single column layout; cards stack vertically.
  \item \textbf{Tablet (\texttt{sm:}):} Sensor cards 2-column grid.
  \item \textbf{Laptop (\texttt{lg:}):} 3-column grid for sensor cards; 3:1 chart+sidebar split.
  \item \textbf{Desktop (\texttt{xl:}):} 5-column sensor card grid; full analytics 2-column layout.
\end{itemize}
